\chapter{Homework 7 (due 10/28)}

\subsection{(Bonus) Lippmann-Schwinger equation}

In this problem, we delve into the derivation of Lippmann-Schwinger equation, which was briefly mentioned in our lectures. Consider a Hamiltonian 
\begin{align}
\hat H=\hat H_0+\hat V
\end{align}
where $\hat H_0$ is the non-interacting Hamiltonian which is exactly solvable, e.g. the free particle Hamiltonian $\hat p^2/2m$ in lectures. $\hat V$ represents an interaction term, e.g. the potential term $V(\hat x)$ for a quantum particle in free space, whose presence makes the full Hamiltonian $\hat H$ hard to solve. 

Next we assume that $\hat H$ and $\hat H_0$ have the same energy spectrum. This is exactly the case for a localized potential
\begin{align}
\lim_{|x|\rightarrow\infty}V(x)=0.
\end{align}
that we have discussed in lectures: it is obvious from the asymptotic behavior of the Schrodinger equation at $|x|\rightarrow\infty$. 

Now we label the orthonormal diagonal basis of non-interacting Hamiltonian $\hat H_0$ to be $\{\dket{a}\}$, satisfying
\begin{align}
\hat H_0\dket{a}=E_a\dket{a},~~~\expval{a|b}=\delta_{a,b}.
\end{align}
We further label the eigenstate of full Hamiltonian $\hat H$ with energy $E_a$ as $\dket{\psi^\pm_{a}}$, satisfying the stationary Schrodinger equation
\begin{align}
\hat H\dket{\psi^\pm_a}=(E_a\pm\imth\epsilon)\dket{\psi^\pm_a}
\end{align}
where $\epsilon\rightarrow0^+$ is an infinitely small positive number. 

(1) \textbf{Bonus} Show that 
\begin{align}
\dket{\psi^\pm_a}=\alpha\dket{a}+(E_a-\hat H_0\pm\imth\epsilon)^{-1}\hat V\dket{\psi^\pm_a}
\end{align}
where $\alpha$ is a complex coefficient.

Comment: we can always absorb the coefficient $\alpha$ into the un-normalized eigenstate $\dket{\psi^\pm_a}$, which leads to the familiar form of Lippmann-Schwinger equation:
\begin{align}
\dket{\psi^\pm_a}=\dket{a}+\frac1{E_a-\hat H_0\pm\imth\epsilon}\hat V\dket{\psi^\pm_a}
\end{align}
But it is sometimes important to keep the coefficient $\alpha$ in mind, as we will see in the next problem when we are dealing with piecewise continuous wavefunctions.\\

(2) \textbf{Bonus} Show that $\{\dket{\psi^+_a}\}$ forms another set of complete orthonormal basis of the Hilbert space. In other words, show that
\begin{align}
\expval{\psi^+_a|\psi^+_b}=\delta_{a,b}=\expval{a|b}
\end{align}

Comment: we need to use the Gellmann-Low theorem for this problem, which we mentioend during lectures. Similarly one can show that $\{\dket{\psi^-_a}\}$ is also a set of complete orthonormal basis. 

\subsection{(Bonus) Solving Lippmann-Schwinger eq. for a delta function potential}

Next, we will use the Lippmann-Schwinger equation to solve the eigenstate wavefunctions for scattering states with a delta function potential:
\begin{align}
\hat H=\hat H_0+\hat V,~~~\hat H_0=\frac{\hat p^2}{2m},~~~\hat V=V(\hat x)
\end{align}
Recall the Lippmann-Schwinger equation:
\begin{align}\label{hw7:lippmann-schwinger eq}
\psi^{\pm}_k(x)=\alpha\expval{x|k}+\int\text{d}yG_0^{\pm}(E|x,y)V(y)\psi^{\pm}_k(y)
\end{align}
where $\alpha$ is an arbitrary coefficient to be determined, and $\dket{k}$ is a plane wave 
\begin{align}
\expval{x|k}=\frac1{\sqrt{2\pi}}e^{\imth kx}
\end{align}
with wavevector $k=\sqrt{2mE}/\hbar>0$, with a positive energy $E>0$ for the free particle Hamiltonian $\hat H_0$. The retarded ($G^+$) and advanced ($G^-$) Green's functions are defined as 
\begin{align}\label{hw7:Green's function}
G_0^{\pm}(E|x,y)=\int\text{d}k\expval{x|k}\frac{1}{E-\hbar^2k^2/2m\pm\imth\epsilon}\expval{k|y}=\int\frac{\text{d}k}{2\pi}\frac{e^{\imth k(x-y)}}{E-\hbar^2k^2/2m\pm\imth\epsilon}
\end{align}
where $\epsilon\rightarrow0^+$ is an infinitely small positive number. 

(1) \textbf{Bonus} Show the following relation holds for an arbitrary potential $V(\hat x)$
\begin{align}
\big[\psi^+_k(x)\big]^\ast=\psi^-_{-k}(x)
\end{align}
using the Lippmann-Schwinger equation (\ref{hw7:lippmann-schwinger eq}). This relation is a consequence of the time reversal symmetry $\hat H^\ast=\hat H$.

(2) \textbf{Bonus} Using the contour integral in complex analysis, compute the retarded Green's function $G_0^+(E|x,y)$ in (\ref{hw7:Green's function}). Explain why it corresponds to an outgoing wave. 

Hint: it should have different behaviors in $x<y$ and $x>y$ regions. 

(3) \textbf{Bonus} Similarly compute the advanced Green's function $G_0^-(E|x,y)$ in (\ref{hw7:Green's function}). Explain why it corresponds to an incoming wave. 

(4) \textbf{Bonus} Now we turn to the delta function potential $V(x)=V_0\delta(x)$. We will focus on only the outgoing states $\dket{\psi^+_k}$. Write down the Lippmann-Schwinger equation (\ref{hw7:lippmann-schwinger eq}) for $x>0$ and $x<0$ separately. Note that the coefficient $\alpha$ can be different for the two regions. Solve the Lippmann-Schwinger equation to obtain the transmission and reflection coefficents $t$ and $r$. They should recover the exact solution we derived in lectures. 

\subsection{S matrix and the phase shift}

In lectures we have mentioned that the $S$ matrix defined as the change of basis between two sets of complete orthonormal basis $\{\dket{\psi^+_a}\}$ and $\{\dket{\psi_a^-}\}$ (see problem 1.(2)~)
\begin{align}
S_{a,b}=\expval{\psi_a^-|\psi_b^+}=\expval{a|U_I(+\infty,-\infty)|b}
\end{align}
is a unitary matrix, satisfying $\hat S^\dagger\hat S=\hat 1$. A unitary matrix can always be diagonalized, and its eigenvalues are phase factors $e^{2\imth\delta}$ known as phase shifts in scattering problems. We will learn more about the phase shift in this problem. 

In our lectures, we have derived the asymptotic form of exact eigenstates for an arbitrary localized potential $V(x)$ satisfying $\lim_{|x|\rightarrow\infty}V(x)=0$:
\begin{align}
\psi^+_k(x)=\begin{cases}
    t e^{\imth kx},~~~&x\rightarrow+\infty,\\
    e^{\imth kx}+re^{-\imth kx},~~~&x\rightarrow-\infty.
\end{cases}\\
\psi^+_{-k}(x)=\begin{cases}
    e^{-\imth kx}+r^\prime e^{\imth kx},~~~&x\rightarrow+\infty,\\
    t^\prime e^{-\imth kx},~~~&x\rightarrow-\infty.
\end{cases}
\end{align}
and 
\begin{align}
\psi^-_k(x)=\big[\psi^+_{-k}(x)\big]^\ast,~~~\psi^-_{-k}(x)=\big[\psi^+_{k}(x)\big]^\ast
\end{align}

(1) Using the continuity of probability flux, show that
\begin{align}\label{hw7:continuity}
|t|^2+|r|^2=1=|t^\prime|^2+|r^\prime|^2.
\end{align}

(2) Using the fact that $(t\psi^+_{-k}-r^\prime\psi_k^+)^\ast\propto\psi_k^+$ and relation (\ref{hw7:continuity}), show that
\begin{align}
t^\prime=t,~~~r^\prime=-r^\ast t/t^\ast.
\end{align}

(3) Using the relations proved in parts (1) and (2), show that 
\begin{align}
\expval{\psi_k^+|\psi^+_{-k}}=\expval{\psi_k^-|\psi^-_{-k}}=0.
\end{align}

Comment: this verifies that either $\{\dket{\psi_k^+},\dket{\psi_{-k}^+}\}$ or $\{\dket{\psi_k^-},\dket{\psi_{-k}^-}\}$ forms a complete orthogonal basis of the Hilbert space. 

(4) \textbf{Bonus} In the subspace with energy $E=\hbar^2k^2/2m>0$, the $\hat S$ matrix is given by
\begin{align}
\hat S(k)=\bpm\expval{\psi^-_{k}|\psi^+_{k}}&\expval{\psi^-_{k}|\psi^+_{-k}}\\ \expval{\psi^-_{-k}|\psi^+_{k}}&\expval{\psi^-_{-k}|\psi^+_{-k}}\epm
\end{align}
Diagonalize it to obtain its eigenvalues $e^{2\imth\delta_+}$ and $e^{2\imth\delta_-}$ as a function of $r$ and $t$. 

Hint: since $|t|^2+|r|^2=1$ we can parameterize them as
\begin{align}
t=\cos\theta e^{\imth\phi_t},~~~r=\sin\theta e^{\imth\phi_r},~~~0\leq\theta\leq\pi/2.
\end{align}

Comment: the diagonal basis of $\hat S(k)$ are inversion eigenstates with inversion eigenvalue $+1$ and $-1$ respectively, if the potential is inversion symmetric i.e. $V(-x)=V(x)$. 

\subsection{Tunneling through a finite potential barrier}

Consider a finite potential barrier
\begin{align}
V(x)=\begin{cases}
    0,~~~&|x|>a/2,\\
    U>0,~~~&|x|<a/2.
\end{cases}
\end{align}
where $a>0$ is the width of the potential barrier. 

(1) Consider an eigenstate, the in $\dket{\psi_k^+}$ with energy $0<E<U$. Write down the form of its eigenstate wavefunction $\psi_k^+(x)$, without computing the coefficients. 

(2) Write down the boundary conditions satisfied by $\psi_k^+(x)$ at $x=\pm a/2$.

(3) Obtain the reflection and transmission coefficients $r$ and $t$ as functions of $E,U,a$. 

(4) Show that in the low energy limit of 
\begin{align}
U-E\gg\frac{\hbar^2}{ma^2}
\end{align}
the transmission probability decays exponentially with $\sqrt{U-E}$
\begin{align}
T\equiv|t|^2\sim e^{-C\sqrt{U-E}}
\end{align}
Compute the coefficient $C$. 

(5) Now we consider $U<0$ case, i.e. a finite potential well which we discussed in details in our lectures. In this case, the even-parity bound state solution is determined by the following equation
\begin{align}
    k'\tan(k'a/2)=\sqrt{\frac{2m|U|}{\hbar^2}-(k')^2},~~~E=\frac{(\hbar k')^2}{2m}-|U|<0.
\end{align}
Show that in the limits of
\begin{align}
    a\rightarrow0^+,~~~|U|\rightarrow\infty,~~~a|U|=|D_0|=\text{const.}
\end{align}
the finite potential well is reduced to an attractive delta function $V(x)=-|D_0|\delta(x)$, and show the bound state energy is given by
\begin{align}
    E=-\frac{m|D_0|^2}{2\hbar^2}
\end{align}

\subsection{A step potential}

Consider the following step potential 
\begin{align}
V(x)=\begin{cases}
    0,~~~&x<0,\\
    U>0,~~~&x>0
\end{cases}
\end{align}

(1) Compute the wavefunction for an eigenstate with energy $E>U$.

(2) Compute the wavefunction for an eigenstate with energy $0<E<U$. 